% \vspace{-0.15in}
\section{Experiments}
\label{sec:experiments}
\begin{table*}[t]
    % Comparison averaged across 4, 8, 16, 32, and 64 input views.
% ETH3D and DTU are averaged across 4, 8, 16, and 32 input views.
\vspace{-0.15in}
\centering
\caption{Quantitative comparison across seven benchmark datasets, averaged over \textbf{4--64 input views}. The top-3 results are highlighted as \colorbox{rankfirst}{first}, \colorbox{ranksecond}{second}, and \colorbox{rankthird}{third}.}
\vspace{0.03in}
\label{tab:main_comparison_4_64view}
\setlength{\tabcolsep}{1pt}
\renewcommand{\arraystretch}{1.28}
\small
\makebox[\textwidth][c]{%
\resizebox{\textwidth}{!}{%
\begin{tabular}{clcccccccccccccc}
    \toprule
    \multirow{2}{*}{\textbf{Input}}
    & \multirow{2}{*}{\textbf{Method}}
    & \multicolumn{2}{c}{\textbf{ScanNet}}
    & \multicolumn{2}{c}{\textbf{NRGBD}}
    & \multicolumn{2}{c}{\textbf{DL3DV}}
    & \multicolumn{2}{c}{\textbf{Tanks \& Temples}}
    & \multicolumn{2}{c}{\textbf{Mip-NeRF 360}}
    & \multicolumn{2}{c}{\textbf{ETH3D}}
    & \multicolumn{2}{c}{\textbf{DTU}} \\
    \cmidrule(lr){3-4}
    \cmidrule(lr){5-6}
    \cmidrule(lr){7-8}
    \cmidrule(lr){9-10}
    \cmidrule(lr){11-12}
    \cmidrule(lr){13-14}
    \cmidrule(lr){15-16}
    & & DISTS$\downarrow$ & LPIPS$\downarrow$
    & DISTS$\downarrow$ & LPIPS$\downarrow$
    & DISTS$\downarrow$ & LPIPS$\downarrow$
    & DISTS$\downarrow$ & LPIPS$\downarrow$
    & DISTS$\downarrow$ & LPIPS$\downarrow$
    & DISTS$\downarrow$ & LPIPS$\downarrow$
    & DISTS$\downarrow$ & LPIPS$\downarrow$ \\
    \midrule
    \multirow{6}{*}{\rotatebox[origin=c]{90}{w/o Pose}}
    & AnySplat
    & 0.278 & 0.414
    & 0.236 & 0.352
    & 0.371 & 0.327
    & 0.593 & 0.737
    & 0.538 & 0.787
    & 0.278 & 0.154
    & 0.520 & 0.810 \\
    & YoNoSplat
    & 0.265 & \third{0.236}
    & 0.285 & 0.240
    & 0.375 & \third{0.219}
    & 0.424 & \second{0.338}
    & 0.405 & 0.586
    & 0.321 & \second{0.033}
    & 0.484 & 0.729 \\
    & DA3-GS
    & 0.286 & 0.305
    & \second{0.233} & \third{0.217}
    & \third{0.337} & 0.228
    & \third{0.414} & 0.380
    & \third{0.371} & \third{0.565}
    & \third{0.228} & 0.087
    & \third{0.365} & \second{0.582} \\
    & Mirror2-GS
    & \second{0.193} & \second{0.222}
    & \third{0.235} & 0.237
    & \second{0.246} & \second{0.188}
    & \second{0.306} & \third{0.351}
    & \second{0.283} & \second{0.528}
    & \second{0.205} & 0.070
    & \second{0.347} & \third{0.617} \\
    & QuerySplat
    & \third{0.236} & 0.242
    & 0.261 & \second{0.216}
    & 0.343 & 0.291
    & 0.425 & 0.393
    & 0.412 & 0.652
    & 0.290 & \first{0.031}
    & 0.424 & 0.676 \\
    & \textbf{Ours}
    & \first{0.184} & \first{0.207}
    & \first{0.195} & \first{0.183}
    & \first{0.221} & \first{0.156}
    & \first{0.252} & \first{0.299}
    & \first{0.233} & \first{0.463}
    & \first{0.147} & \third{0.049}
    & \first{0.229} & \first{0.483} \\
    \midrule
    \multirow{7}{*}{\rotatebox[origin=c]{90}{w/ Pose}}
    & DepthSplat
    & 0.288 & 0.278
    & 0.323 & 0.272
    & 0.329 & 0.191
    & 0.397 & 0.353
    & 0.561 & 0.656
    & 0.445 & 0.068
    & 0.370 & 0.598 \\
    & TokenGS
    & 0.369 & 0.272
    & 0.333 & 0.277
    & 0.375 & 0.211
    & 0.515 & 0.346
    & 0.502 & 0.716
    & 0.400 & \third{0.042}
    & 0.593 & 0.828 \\
    & tttLRM
    & 0.292 & 0.300
    & 0.320 & 0.292
    & 0.323 & 0.196
    & \third{0.390} & \third{0.342}
    & 0.350 & 0.531
    & 0.225 & 0.063
    & 0.382 & 0.599 \\
    & ReSplat
    & \third{0.272} & \third{0.240}
    & 0.284 & 0.232
    & \third{0.304} & \second{0.159}
    & 0.396 & \first{0.275}
    & \third{0.349} & \third{0.525}
    & \third{0.219} & \first{0.027}
    & 0.441 & 0.664 \\
    & DA3-GS (Pose Cond.)
    & 0.274 & 0.291
    & \third{0.238} & \third{0.227}
    & 0.311 & 0.219
    & 0.425 & 0.368
    & 0.362 & 0.556
    & 0.223 & 0.089
    & \third{0.368} & \third{0.596} \\
    & Mirror2-GS (Pose Cond.)
    & \second{0.195} & \second{0.223}
    & \second{0.231} & \second{0.226}
    & \second{0.239} & \third{0.180}
    & \second{0.304} & 0.346
    & \second{0.272} & \second{0.506}
    & \second{0.201} & 0.062
    & \second{0.315} & \second{0.532} \\
    & \textbf{Ours (Pose Cond.)}
    & \first{0.185} & \first{0.206}
    & \first{0.186} & \first{0.162}
    & \first{0.215} & \first{0.148}
    & \first{0.251} & \second{0.290}
    & \first{0.219} & \first{0.440}
    & \first{0.133} & \second{0.037}
    & \first{0.184} & \first{0.364} \\
    \bottomrule
\end{tabular}%
}%
}
\vspace{-0.05in}

    % % Comparison across seven benchmark datasets.
% Requires \usepackage{graphicx}, which is already loaded in the main file.
    \centering
    \caption{Quantitative comparison across seven benchmark datasets using \textbf{32 input views}. The top-3 results are highlighted as \colorbox{rankfirst}{first}, \colorbox{ranksecond}{second}, and \colorbox{rankthird}{third}.}
    % \vspace{0.05in}
    \label{tab:main_comparison_32view}
    \setlength{\tabcolsep}{1pt}
    \renewcommand{\arraystretch}{1.2}
    \small
    \makebox[\textwidth][c]{%
    \resizebox{\textwidth}{!}{%
    \begin{tabular}{clcccccccccccccc}
        \toprule
        \multirow{2}{*}{\textbf{Input}}
        & \multirow{2}{*}{\textbf{Method}}
        & \multicolumn{2}{c}{\textbf{ScanNet}}
        & \multicolumn{2}{c}{\textbf{NRGBD}}
        & \multicolumn{2}{c}{\textbf{DL3DV}}
        & \multicolumn{2}{c}{\textbf{Tanks \& Temples}}
        & \multicolumn{2}{c}{\textbf{Mip-NeRF 360}}
        & \multicolumn{2}{c}{\textbf{ETH3D}}
        & \multicolumn{2}{c}{\textbf{DTU}} \\
        \cmidrule(lr){3-4}
        \cmidrule(lr){5-6}
        \cmidrule(lr){7-8}
        \cmidrule(lr){9-10}
        \cmidrule(lr){11-12}
        \cmidrule(lr){13-14}
        \cmidrule(lr){15-16}
        & & DISTS$\downarrow$ & LPIPS$\downarrow$
        & DISTS$\downarrow$ & LPIPS$\downarrow$
        & DISTS$\downarrow$ & LPIPS$\downarrow$
        & DISTS$\downarrow$ & LPIPS$\downarrow$
        & DISTS$\downarrow$ & LPIPS$\downarrow$
        & DISTS$\downarrow$ & LPIPS$\downarrow$
        & DISTS$\downarrow$ & LPIPS$\downarrow$ \\
        \midrule
        \multirow{6}{*}{\rotatebox[origin=c]{90}{w/o Pose}}
        & AnySplat
        & 0.300 & 0.430
        & \second{0.232} & 0.342
        & 0.380 & 0.338
        & 0.655 & 0.811
        & 0.618 & 0.833
        & 0.241 & 0.106
        & 0.577 & 0.805 \\
        & YoNoSplat
        & 0.268 & \third{0.238}
        & 0.296 & 0.253
        & 0.382 & \third{0.225}
        & 0.430 & \second{0.342}
        & 0.400 & 0.581
        & 0.316 & \first{0.039}
        & 0.497 & 0.706 \\
        & DA3-GS
        & 0.270 & 0.282
        & 0.249 & \third{0.233}
        & 0.355 & 0.234
        & \third{0.379} & 0.358
        & \third{0.352} & \third{0.557}
        & \third{0.211} & 0.091
        & \second{0.314} & \second{0.462} \\
        & Mirror2-GS
        & \second{0.192} & \second{0.229}
        & \third{0.238} & 0.247
        & \second{0.249} & \second{0.195}
        & \second{0.312} & \third{0.352}
        & \second{0.288} & \second{0.544}
        & \second{0.204} & 0.082
        & \third{0.371} & \third{0.623} \\
        & QuerySplat
        & \third{0.231} & 0.253
        & 0.263 & \second{0.223}
        & \third{0.343} & 0.301
        & 0.428 & 0.400
        & 0.408 & 0.658
        & 0.286 & \second{0.043}
        & 0.403 & 0.627 \\
        & \textbf{Ours}
        & \first{0.181} & \first{0.205}
        & \first{0.199} & \first{0.182}
        & \first{0.223} & \first{0.154}
        & \first{0.253} & \first{0.300}
        & \first{0.231} & \first{0.460}
        & \first{0.161} & \third{0.071}
        & \first{0.234} & \first{0.455} \\
        \midrule
        \multirow{7}{*}{\rotatebox[origin=c]{90}{w/ Pose}}
        & DepthSplat
        & \third{0.238} & 0.265
        & 0.313 & 0.277
        & \third{0.316} & 0.189
        & \third{0.340} & \third{0.341}
        & 0.530 & 0.646
        & 0.232 & \first{0.035}
        & 0.373 & 0.560 \\
        & TokenGS
        & 0.374 & 0.277
        & 0.332 & 0.282
        & 0.378 & 0.213
        & 0.534 & 0.354
        & 0.519 & 0.740
        & 0.347 & \second{0.043}
        & 0.554 & 0.772 \\
        & tttLRM
        & 0.281 & 0.286
        & 0.292 & 0.250
        & \third{0.316} & 0.190
        & 0.386 & 0.346
        & 0.334 & 0.525
        & \third{0.199} & 0.066
        & 0.349 & 0.528 \\
        & ReSplat
        & 0.271 & \third{0.245}
        & 0.302 & 0.249
        & 0.328 & \second{0.170}
        & 0.368 & \first{0.263}
        & \third{0.323} & \second{0.517}
        & 0.242 & \first{0.035}
        & 0.343 & \third{0.512} \\
        & DA3-GS (Pose Cond.)
        & 0.254 & 0.259
        & \third{0.255} & \third{0.244}
        & 0.336 & 0.235
        & 0.397 & 0.354
        & 0.340 & 0.537
        & 0.202 & 0.088
        & \third{0.324} & \second{0.473} \\
        & Mirror2-GS (Pose Cond.)
        & \second{0.194} & \second{0.230}
        & \second{0.237} & \second{0.241}
        & \second{0.241} & \third{0.185}
        & \second{0.309} & 0.346
        & \second{0.277} & \third{0.523}
        & \second{0.193} & 0.070
        & \second{0.322} & 0.521 \\
        & \textbf{Ours (Pose Cond.)}
        & \first{0.184} & \first{0.212}
        & \first{0.186} & \first{0.161}
        & \first{0.213} & \first{0.147}
        & \first{0.254} & \second{0.303}
        & \first{0.219} & \first{0.452}
        & \first{0.139} & \third{0.057}
        & \first{0.185} & \first{0.357} \\
        \bottomrule
    \end{tabular}%
    }%
    }
    \vspace{-0.15in}
\end{table*}

% \vspace{-0.1in}
\subsection{Evaluation Benchmark}
Existing benchmarks are typically concentrated on a small set of datasets and evaluate interpolated target views, making it difficult to assess cross-domain and viewpoint generalization. To provide a more comprehensive evaluation, we design a new benchmark with the following properties:

% \vspace{-0.02in}

\textbf{Cross-domain and cross-scale coverage.}
We evaluate on \textbf{seven datasets}: ScanNet~\citep{dai2017scannet}, NRGBD~\citep{azinovic2022neural}, DL3DV~\citep{ling2024dl3dv}, Tanks and Temples~\citep{knapitsch2017tanks}, Mip-NeRF 360~\citep{barron2022mip}, ETH3D~\citep{schops2017multi}, and DTU~\citep{jensen2014large}. These datasets span a wide range of physical scene scales, providing a substantially broader test of cross-domain and cross-scale generalization.

% \vspace{-0.02in}

\textbf{Controlled extrapolation and observation sparsity.}
We sample context views from a compact trajectory span and select target
views using an occlusion-aware symmetric co-visibility criterion.
The input--target overlap is controlled within \textbf{0.2--0.6}, with
lower overlap corresponding to more challenging extrapolation.
We further evaluate \textbf{1/4/8/16/32/64 input views} under a fixed
context span to isolate the effect of observation density.
More details are provided in the appendix.

\begin{figure*}[t] 
    \centering
    \includegraphics[width=\textwidth]{figures/exp/main_compare.jpg}
    \vspace{-0.2in}
    \caption{Qualitative comparisons across diverse scenes and viewpoints with multi-view input.}
    \label{fig:qualitative}
    \vspace{-0.2in}
\end{figure*}


\textbf{Multi-resolution.}
Our benchmark covers a broad resolution range, from approximately \textbf{480p to 1K}, enabling evaluation under both low- and high-resolution settings.

\textbf{Compared methods.}
We evaluate two categories methods according to whether camera poses are provided. In the \emph{w/o Pose} setting, methods reconstruct the scene directly from RGB observations, including AnySplat~\citep{jiang2025anysplat}, YoNoSplat~\citep{ye2026yonosplat}, DA3-GS~\citep{lin2025depth}, Mirror2-GS~\citep{team2026hy}, and QuerySplat~\citep{li2026querysplat}.

\begin{wraptable}{r}{0.50\textwidth}
    \vspace{-0.15in}
    % Monocular comparison across seven benchmark datasets.
    \centering
    \caption{Quantitative comparison of feed-forward 3DGS methods using a \textbf{single input view}. The top-3 results are highlighted as \colorbox{rankfirst}{first}, \colorbox{ranksecond}{second}, and \colorbox{rankthird}{third}.}
    \vspace{0.01in}
    \label{tab:main_comparison_monocular}
    \setlength{\tabcolsep}{4pt}
    \renewcommand{\arraystretch}{1.15}
    \small
    \makebox[\linewidth][c]{%
    \resizebox{\linewidth}{!}{%
    \begin{tabular}{lcccccccccccccc}
        \toprule
        \multirow{2}{*}{\textbf{Method}}
        & \multicolumn{2}{c}{\textbf{ScanNet}}
        & \multicolumn{2}{c}{\textbf{NRGBD}}
        & \multicolumn{2}{c}{\textbf{DL3DV}}
        & \multicolumn{2}{c}{\textbf{Tanks \& Temples}}
        & \multicolumn{2}{c}{\textbf{Mip-NeRF 360}}
        & \multicolumn{2}{c}{\textbf{ETH3D}}
        & \multicolumn{2}{c}{\textbf{DTU}} \\
        \cmidrule(lr){2-3}
        \cmidrule(lr){4-5}
        \cmidrule(lr){6-7}
        \cmidrule(lr){8-9}
        \cmidrule(lr){10-11}
        \cmidrule(lr){12-13}
        \cmidrule(lr){14-15}
        & DISTS$\downarrow$ & LPIPS$\downarrow$
        & DISTS$\downarrow$ & LPIPS$\downarrow$
        & DISTS$\downarrow$ & LPIPS$\downarrow$
        & DISTS$\downarrow$ & LPIPS$\downarrow$
        & DISTS$\downarrow$ & LPIPS$\downarrow$
        & DISTS$\downarrow$ & LPIPS$\downarrow$
        & DISTS$\downarrow$ & LPIPS$\downarrow$ \\
        \midrule
        SHARP
        & \first{0.223} & \first{0.244}
        & \first{0.238} & \first{0.174}
        & \second{0.335} & \second{0.280}
        & \second{0.374} & \second{0.470}
        & \first{0.328} & \first{0.505}
        & \second{0.200} & \second{0.133}
        & \second{0.433} & \second{0.724} \\
        InfiniSplat
        & \third{0.266} & \third{0.280}
        & \third{0.336} & \third{0.345}
        & \third{0.527} & \third{0.596}
        & \third{0.452} & \third{0.531}
        & \third{0.505} & \third{0.654}
        & \third{0.367} & \third{0.376}
        & \third{0.562} & \third{0.754} \\
        \textbf{Ours}
        & \second{0.238} & \second{0.274}
        & \second{0.265} & \second{0.284}
        & \first{0.279} & \first{0.220}
        & \first{0.342} & \first{0.430}
        & \second{0.410} & \second{0.599}
        & \first{0.195} & \first{0.104}
        & \first{0.327} & \first{0.657} \\
        \bottomrule
    \end{tabular}%
    }%
    }
    \vspace{-0.1in}
    \par\vspace{0.1in}
    \centering
    \includegraphics[width=\linewidth]{figures/exp/main_compare_mono.jpg}
    \vspace{-0.23in}
    \captionof{figure}{Qualitative comparisons across diverse scenes and viewpoints with monocular input.}
    \label{fig:qualitative_mono}
    \vspace{-0.5in}
\end{wraptable}

In the \emph{w/ Pose} setting, camera poses are additionally provided for model input. We compare against DepthSplat~\citep{xu2025depthsplat}, TokenGS~\citep{ren2026tokengs}, tttLRM~\citep{wang2026tttlrm}, ReSplat~\citep{xu2025resplat}, as well as pose-conditioned variants of DA3-GS~\citep{lin2025depth} and Mirror2-GS~\citep{team2026hy}.
Additionally, we compare against SHARP~\citep{mescheder2026sharp} and InfiniSplat~\citep{wang2026infinisplat}, which target for monocular view synthesis.

% \vspace{-0.1in}
\subsection{Quantitative Comparisons}
% \vspace{-0.07in}

We report the average metrics over \textbf{4/8/16/32/64 input views} in Table~\ref{tab:main_comparison_4_64view}, specific results under different input-view settings are provided in the appendix. Following SHARP~\citep{mescheder2026sharp}, we primarily report \textit{LPIPS} and \textit{DISTS} in the main paper, as they better reflect rendering quality and higher-frequency details. We also report the results in the monocular setting in Table~\ref{tab:main_comparison_monocular}.
Overall, SplatAnything achieves the strongest performance across most datasets and evaluation settings, demonstrating its superior generalization capabilities.

% \vspace{-0.1in}
\subsection{Qualitative Comparisons}
% \vspace{-0.1in}

Figs.~\ref{fig:qualitative} and~\ref{fig:qualitative_mono} present qualitative comparisons across diverse scenes and extrapolated viewpoints under multi-view and monocular input settings, covering object-centric captures, indoor environments, and large-scale outdoor scenes. \textbf{SplatAnything} demonstrates a significant advantage in rendering quality: it generalizes robustly across a wide range of scene types and spatial scales, maintains geometric consistency under challenging large viewpoint
extrapolation, and produces sharp, detailed renderings that faithfully recover fine local structures. More results are provided in the appendix.
% \vspace{-0.08in}
\subsection{Efficiency Analysis}
% \vspace{-0.08in}
We evaluate the inference efficiency of \textbf{SplatAnything} and other methods under different input-view settings. Detailed timing results are provided in the appendix.

% \vspace{-0.08in}
\subsection{Ablation Studies}
% \vspace{-0.06in}
\paragraph{Effect of Gaussian spacing normalization.}
We ablate the normalization space for Gaussian prediction in
Table~\ref{tab:ablation_normalization_space}, comparing RDNorm,
RDNorm with scale calibration, and GSNorm under the same architecture
and training settings.
For RDNorm with scale calibration, Gaussian prediction remains in the
RDNorm space, while the predicted Gaussian scales are uniformly rescaled
to match the scene's median Gaussian spacing. Both RDNorm variants remain consistently inferior to GSNorm.
Qualitatively, RDNorm tends to produce holes in small object-centric
scenes and over-smoothed results in large outdoor scenes
(Fig.~\ref{fig:ablation}).
Scale calibration alleviates the holes on object-centric scenes, but
often turns them into over-smoothed reconstructions, while outdoor scenes
remain blurry.
These results indicate that post-prediction scale correction cannot
resolve the ambiguity in scale supervision introduced by RDNorm.
By canonicalizing the representative Gaussian spacing scale before Gaussian
prediction, GSNorm provides a better-aligned prediction space and yields
more complete and sharper reconstructions across diverse scenes. 
Additional results are provided in the appendix.

% \vspace{-0.1in}
\paragraph{Training design choices}
We also analyze several architecture and training design choices, including
fine-tuning the pretrained geometry backbone, fixing Gaussian centers,
scaling the amount of training data, and using mixed sampling strategies
for Gaussian supervision. Please refer to the appendix for more analysis and results.

% Ablation study on the choice of normalization space.
% Replace ``--'' with the corresponding experimental results.
\vspace{-0.04in}
\begin{table*}[t]
    \centering
    \caption{Normalization-space ablation across different numbers of input views.}
    \vspace{0.02in}
    \label{tab:ablation_normalization_space}
    \setlength{\tabcolsep}{4pt}
    \renewcommand{\arraystretch}{1.05}
    \small
\makebox[\textwidth][c]{%
\resizebox{\textwidth}{!}{%
    \begin{tabular}{lcccccccccccc}
        \toprule
        \multirow{2}{*}{\textbf{Norm. Space}}
        & \multicolumn{2}{c}{\textbf{Num. Views=1}}
        & \multicolumn{2}{c}{\textbf{Num. Views=4}}
        & \multicolumn{2}{c}{\textbf{Num. Views=8}}
        & \multicolumn{2}{c}{\textbf{Num. Views=16}}
        & \multicolumn{2}{c}{\textbf{Num. Views=32}}
        & \multicolumn{2}{c}{\textbf{Num. Views=64}} \\
        \cmidrule(lr){2-3}
        \cmidrule(lr){4-5}
        \cmidrule(lr){6-7}
        \cmidrule(lr){8-9}
        \cmidrule(lr){10-11}
        \cmidrule(lr){12-13}
        & DISTS$\downarrow$ & LPIPS$\downarrow$
        & DISTS$\downarrow$ & LPIPS$\downarrow$
        & DISTS$\downarrow$ & LPIPS$\downarrow$
        & DISTS$\downarrow$ & LPIPS$\downarrow$
        & DISTS$\downarrow$ & LPIPS$\downarrow$
        & DISTS$\downarrow$ & LPIPS$\downarrow$ \\
        \midrule
        Radial Distance Norm.
        & 0.308 & 0.394
        & 0.223 & 0.278
        & 0.223 & 0.272
        & 0.218 & 0.268
        & 0.225 & 0.277
        & 0.224 & 0.250 \\
        Radial Distance Norm + Scale Calibration.
        & 0.323 & 0.400
        & 0.249 & 0.309
        & 0.252 & 0.308
        & 0.245 & 0.300
        & 0.253 & 0.305
        & 0.234 & 0.265 \\
        \textbf{Gaussian Spacing Norm. (Ours)}
        & \textbf{0.294} & \textbf{0.367}
        & \textbf{0.198} & \textbf{0.245}
        & \textbf{0.196} & \textbf{0.237}
        & \textbf{0.190} & \textbf{0.232}
        & \textbf{0.197} & \textbf{0.241}
        & \textbf{0.217} & \textbf{0.229} \\
        \bottomrule
    \end{tabular}%
}%
}
\end{table*}
% \vspace{-0.2in}
\begin{figure*}[t] 
    \centering
    \vspace{-0.1in}
    \includegraphics[width=\linewidth]{figures/exp/ablation.jpg}
    \vspace{-0.25in}
    \caption{Ablation studies on the effect of GSNorm across different scene types.}
    \label{fig:ablation}
    \vspace{-0.12in}
\end{figure*}

% Quantitative and qualitative results are shown in Tab.~\ref{tab:ablation_training_design_choices} and Fig.~\ref{fig:ablation_training_design_choices}. We find that finetuning the pretrained geometry backbone brings improvements. Fixing the gaussian centers ensures more correct geometry across viewpoint changes, as RGB-projection loss tends to overfit the geometry to the input views. Also, using mixed sampling strategies for training feed-forward Gaussian Splatting can learn more reaonable Gaussian attributes, especially under the extrapolation viewpoint shifts.


